\begin{abcliste}
\abc Die Formelsammlung (FoTaBe S.~210) gibt die mittlere Entfernung $r_\mathrm{EM} = \rEMO$ und die siderische Umlaufzeit $T = \TMO = \TM$. Für eine Kreisbahn ist die Radialbeschleunigung
\begin{align}
 a_R &= \frac{v^2}{r_\mathrm{EM}} = \aRF \\
   &= \frac{4\pi^2\cdot\rEM}{\left(\TM\right)^2} \\
   &= \aR \approx \result{\aRP{-3}{3}}
\end{align}

\abc An der Erdoberfläche erzeugt die Gravitationskraft der Erde die Fallbeschleunigung, im Abstand $r_\mathrm{EM}$ die Radialbeschleunigung des Mondes:
\begin{align}
 m\,g &= G\,\frac{M_\mathrm{E}\,m}{r_\mathrm{E}^2} &
 m_\mathrm{M}\,a_R &= G\,\frac{M_\mathrm{E}\,m_\mathrm{M}}{r_\mathrm{EM}^2}
\end{align}
Dividiert man die beiden Gleichungen $g = G M_\mathrm{E}/r_\mathrm{E}^2$ und $a_R = G M_\mathrm{E}/r_\mathrm{EM}^2$, fallen $G$ und $M_\mathrm{E}$ heraus:
\begin{align}
 \frac{g}{a_R} &= \frac{r_\mathrm{EM}^2}{r_\mathrm{E}^2} = \left(\frac{r_\mathrm{EM}}{r_\mathrm{E}}\right)^2 \quad\square
\end{align}

\abc Mit $g = \gEO$ und dem Resultat von a) ist das gemessene Verhältnis
\begin{align}
 \frac{g}{a_R} &= \qaF = \frac{\gE}{\aR} = \qa \approx \result{\qaP{0}{3}}
\end{align}
Das Gravitationsgesetz sagt mit $r_\mathrm{E} = \RErO$ (FoTaBe S.~209) voraus:
\begin{align}
 \left(\frac{r_\mathrm{EM}}{r_\mathrm{E}}\right)^2 &= \qbF = \left(\frac{\rEM}{\REr}\right)^2 = \qb \approx \result{\qbP{0}{3}}
\end{align}
Der vorausgesagte Wert ist nur um $\dqP{0}{1}$ grösser als der gemessene. Die Beschleunigung des Mondes ist also genau so viel kleiner als die Fallbeschleunigung des Apfels, wie es eine mit $1/r^2$ abnehmende Erdanziehung verlangt: Apfel und Mond werden von derselben Kraft angezogen. Die kleine Abweichung kommt u.\,a. daher, dass der Mond nicht um den Erdmittelpunkt, sondern um den gemeinsamen Schwerpunkt kreist, seine Bahn leicht elliptisch ist und $g$ den Einfluss der Erdrotation enthält.
\end{abcliste}
